% Folder 136, batch 08 — archivists' pages 141–160.
% Pass: Opus 5.5 (claude-opus-5-5), 2026-10-03 — first pass, unchecked against the pages by a human.
%
% Skipped (no \page): 141, 153, 155, 157 (typed pages of an SGA 7 exposé IX
% typescript, nothing in his hand bearing on these notes), 142 (computer
% listing), 160 (blank). Page 159 is a typescript leaf of the same kind; only
% his handwritten lines on it are transcribed.
\documentclass[11pt,a4paper]{article}
\input{../preamble/grothendieck}

\folder{136}
\batch{8}
\pages{141}{160}
\foldertitle{Complexe de De Rham à puissance divisée [conférence de 1976 à l’IHÉS] : notes manuscrites (s.d.).}
\dating{[à partir de 1975-1976]}
\watermark{Édition de démonstration}

\begin{document}

\section{Le yoga coalgèbres $\mapsto$ Algèbres}

\page{143}
\note{titre encadré en tête de page : « Le yoga coalgèbres $\mapsto$ Algèbres \ill{} » ; il ouvre une suite continue jusqu'à la p. 151.}

$k$ anneau, on travaille dans le topos $\widehat{C}$, où $C_k$ est la catégorie des $k$-algèbres. Le foncteur oubli
\[
  C_k \longrightarrow \mathrm{Ens}, \qquad A \longmapsto A
\]
est un objet de $\widehat{C}_k$ muni d'une structure de $k$-Algèbre, \struck{foncteur vers les $k$-algèbres} noté $\mathcal{O}_k$. Si $M$ est un $k$-module, $V(M)$ (ou $V_k(M)$) \uncertain{désignera} le \struck{foncteur} $\mathcal{O}_k$-Module
\[
  V(M)(A) = \operatorname{Hom}_{k\text{-mod}}(M, A) = \operatorname{Hom}_{A\text{-mod}}(M_A, A).
\]
$V(M)$ dépend fonctoriellement de $M$, de façon contravariante, et on voit que $V$ est un foncteur \uncertain{pl. fid.} des $k$-modules vers les $\mathcal{O}_k$-Modules
\[
  \operatorname{Hom}_k(M, N) \xrightarrow{\ \sim\ } \operatorname{Hom}_{\mathcal{O}_k}(V(N), V(M)).
\]
De façon générale\add{, comme $V(N)$ est représentable par $\operatorname{Sym}^*(N)$,} les morphismes dans $\widehat{C}_k$
\[
  V(N) \longrightarrow \text{\struck{$V(M), \operatorname{Sym}(B)$}}\ F \qquad (\text{où } F \in \operatorname{ob} \widehat{C}_k)
\]
correspondent aux éléments de $F(\operatorname{Sym}^*(N))$ :
\[
  \operatorname{Hom}_{\widehat{C}_k}(V(N), F) \simeq F(\operatorname{Sym}^*(N)),
\]
en particulier si $F = \operatorname{Spec} B$
\[
  \operatorname{Hom}_{\widehat{C}_k}(V(N), \operatorname{Spec} B) \simeq \operatorname{Hom}_{k\text{-alg}}(B, \operatorname{Sym}^*(N)) ;
\]
en faisant $B = \operatorname{Sym}^*(M)$ il vient
\[
  \operatorname{Hom}_{\widehat{C}_k}(V(N), V(M)) \simeq \operatorname{Hom}_{k\text{-mod}}(M, \operatorname{Sym}^* N).
\]
Soit
\[
  u^* : V(N) \longrightarrow V(M)
\]
correspondant à
\[
  u : M \longrightarrow \operatorname{Sym}^*(N).
\]
On constate que \struck{\ill{}}
\[
  u(x) = \sum u_i(x), \qquad u_i(x) \in \operatorname{Sym}^i(N)
\]
(la famille des \struck{\ill{}}
\[
  u_i : M \longrightarrow \operatorname{Sym}^i(N)
\]
\struck{\ill{}} « loc. finie »), alors si on a $\lambda \in \mathcal{O}_k(A) = A$

\page{144}
et $x^* \in V(N)(A) = \operatorname{Hom}_{k\text{-mod}}(N, A)$,
\[
  u^*(x^*) = \widetilde{x^*} \circ u \in \operatorname{Hom}_{k\text{-mod}}(M, A), \qquad M \xrightarrow{\ u\ } \operatorname{Sym}^* N \xrightarrow{\ \widetilde{x^*}\ } A,
\]
où $\widetilde{x^*}$ est l'hom.\ de $k$-alg.\ prolongeant $x^*$. On a donc
\[
  u^*(x^*)(x) = \sum_{i \geqslant 0} \widetilde{x^*}^{\,i}(u_i(x)), \qquad x \in M
\]
($\widetilde{x^*}^{\,i}$ étant la \uncertain{restriction} de $x^*$ à $\operatorname{Sym}^i N$), d'où, si $\lambda \in A$,
\[
  u^*(\lambda x^*)(x) = \sum_{i \geqslant 0} \lambda^i\, \widetilde{x^*}^{\,i}(u_i(x)) = \sum_{i \geqslant 0} \widetilde{x^*}^{\,i}(\lambda^i u_i(x)).
\]
Donc dans la correspondance entre $u^*$ et $u$, \ill{} $A$, \ill{} de \ill{} $x^*$ \ill{} la \struck{multiplication} \add{\uncertain{correspondance}} \struck{\ill{}} \ill{} $\lambda$ dans le $\mathcal{O}_k$-module \struck{\ill{}} \ill{} $N_A$) correspond l'opération
\[
  u = \sum u_i \longmapsto u^\lambda = \sum \lambda^i u_i .
\]
On en conclut que $u^*$ est \uncertain{homogène de degré} $n$ (i.e.\ satisfait, pour $A$, $\lambda$ variables,
\[
  u^* \circ \lambda_{V(N)} = \lambda^n_{V(M)} \circ u^* \text{)}
\]
ssi tous les $u_i$ pour $i \neq n$ sont nuls, i.e.\ ssi $u$ est un hom
\[
  M \longrightarrow \operatorname{Sym}^n N.
\]
[En effet, il faut avoir $\sum (\lambda^i - \lambda^n) u_i = 0$ $\forall A$, $\lambda \in A$, prenant $A = k[T]$ et $\lambda = T$, on doit avoir $\sum (T^i - T^n) u_i(x) = 0$, i.e.\ $\sum (T^i - T^n) u_i(x) = 0$ pour tout $i \neq n$, $x \in M$, $n \neq i$, d'où \struck{$u_i(x)$} $u_i(x) = 0$

\page{145}
pour $i \neq n$, et (comme $x$ arbitraire) $u_i = 0$ pour $i \neq n$.] Donc les \struck{\ill{}} \add{morphismes} homogènes de degré $1$ de $V(N)$ dans $V(M)$ correspondent exactement aux homs de $k$-modules
\[
  M \longrightarrow N
\]
(qui sont par suite nécessairement $\mathcal{O}_k$-linéaires !). Cela prouve en même temps la pleine fidélité.

On détermine de même les morphismes
\[
  \prod_{i \in I} V(N_i) \longrightarrow F \qquad \text{\struck{$V(\textstyle\coprod_i N_i)$}}
\]
compte tenu de l'isom.\ évident
\[
  \prod_{\alpha \in I} V(N_\alpha) \simeq V\Bigl(\bigoplus N_\alpha\Bigr)
\]
\note{double trait vertical en marge gauche des deux lignes qui suivent.}
(de façon générale il est trivial que $V$ transforme $\varinjlim$ quelc.\ en $\varprojlim$). Dans le cas de $F = V(M)$, cela correspond donc aux hom
\[
  M \xrightarrow{\ u\ } \operatorname{Sym}^*\Bigl(\bigoplus N_\alpha\Bigr) = \bigotimes_\alpha \operatorname{Sym}^*(N_\alpha).
\]
Ici il y a lieu de décomposer $M$ suivant les multidegrés $i = (i_\alpha)_{\alpha \in I} \in \mathbb{N}^I$,
\[
  u(x) = \sum_{(i_\alpha)} \underbrace{u_{(i_\alpha)}(x)}_{\in\, \bigotimes_\alpha \operatorname{Sym}^{i_\alpha}(N_\alpha)}
\]

\page{146}
\ill{} on voit que si $\lambda = (\lambda_\alpha) \in \mathcal{O}_k(A)^I$, $(x^*_\alpha) \in \bigl(\prod V(N_\alpha)\bigr)(A)$,
\[
  u^*\bigl((\lambda_\alpha x^*_\alpha)_{\alpha \in I}\bigr) = (u^\lambda)^*(x^*_\alpha)
\]
i.e.
\[
  u^* \prod_\alpha \lambda_{\alpha\, V(N_\alpha)} = (u^\lambda)^* ,
\]
où \struck{$u^\lambda$}
\[
  u^\lambda = \sum_{(i_\alpha)} \Bigl(\prod_\alpha \lambda_\alpha^{i_\alpha}\Bigr) u_{(i_\alpha)} .
\]
Donc, si on se donne $n = (n_\alpha)_{\alpha \in I} \in \mathbb{N}^I$, $u^*$ est multihomogène de \add{multi}degré $n$ ssi \uncertain{c'est} \ill{}
\[
  u : M \longrightarrow \bigotimes_\alpha \operatorname{Sym}^{n_\alpha}(N_\alpha).
\]
En particulier, si tous les $n_\alpha$ sont égaux à $1$, cela signifie que $u^*$ est homogène de degré $1$ en tous les arguments ssi
\[
  u : M \longrightarrow \bigotimes_\alpha N_\alpha
\]
(en identifiant $N_\alpha$ à $\operatorname{Sym}^1 N_\alpha$). Mais alors on constate tout de suite que $u^*$ est \add{\ill{}} multilinéaire (prouvant aussi \struck{\ill{}} \ill{} \ill{} du cas d'$1$ variable (« en fixant les autres »)). Donc les accouplements multil.
\[
  \prod V(N_\alpha) \longrightarrow \text{\struck{\ill{}}}\ V(M)
\]
correspondent aux hom
\[
  M \longrightarrow \bigotimes_\alpha N_\alpha ,
\]
en particulier les accouplements bil.
\[
  V(N) \times V(N') \longrightarrow V(M)
\]
correspondent aux hom
\[
  M \longrightarrow N \otimes N' .
\]

\page{147}
Faisant $M = N = N'$, on trouve qu'une loi de composition bilinéaire sur $V(M)$ correspond à un hom
\[
  \Delta : M \longrightarrow M \otimes M
\]
\note{à droite de la formule : « on dit que $M$ devient une \add{coalgèbre} », la suite de la phrase biffée.}
(\uncertain{cf.} aussi « application diagonale ») ; \struck{on dit que} $\Delta$ est associative, \add{\uncertain{resp.}} \struck{commutative}, \add{\uncertain{resp.}} \struck{unitaire} si la loi de composition sur $V(M)$ l'est. \struck{\ill{}} \uncertain{Cela s'explicite} par les diagrammes bien connus. Pour l'unitarité, on note que $e \in V(M)(k)$ s'identifie à une augmentation ou counité
\[
  M \xrightarrow{\ \varepsilon\ } k
\]
(satisfaisant des conditions qu'on n'écrit pas, qui se définissent de façon unique). Par contre, une augmentation (\uncertain{ou multiplication}) $V(M) \to \mathcal{O}_k = V(k)$ correspond à un hom $k \to M$, i.e.\ à un élément de $M$ (\uncertain{appelé unité}) la condition ici est
\[
  \Delta e = e \otimes e .
\]
Notons que par suite de la formule
\[
  V\bigl(\varinjlim_i M_i\bigr) = \varprojlim_i V(M_i)
\]
(pour \struck{\ill{}} \add{lim.}\ d'indices quelc.), une lim.\ quelc.\ de coalgèbres est une coalgèbre, et les \struck{\ill{}} \add{axiomes} d'ass., d'unit., de comm.

\page{148}
passent à la limite.

Soit \struck{\ill{}} $Z$ un ens. Une graduation de type $Z$ de $M$ est une décomposition en somme directe
\[
  M \simeq \bigoplus_{\alpha \in Z} M_\alpha \qquad ((M_\alpha) \text{ sous-modules de } M)
\]
\struck{\ill{}} s'écrit aussi comme décomposition de $V(M)$ en produit
\[
  V(M) \simeq \prod_{\alpha \in Z} V(M_\alpha)
\]
(notations évidentes). Il est \uncertain{d'ailleurs} commode \add{(\ill{} \ill{} de cela)} \struck{\ill{}} \add{les} \ill{} les \add{$M_*$ gradués}, par $Z$ comme des familles $(M_\alpha)_{\alpha \in Z}$ de $k$-modules, donc $V(M_*)$ comme la famille des $V(M_\alpha)$ $(\alpha \in Z)$.

Si $Z$ est muni d'une structure de groupe abélien, alors la structure dite sur $M$ revient à faire opérer le groupe diagonalisable de $\widehat{C}_k$
\[
  D(Z) = G : A \longmapsto \operatorname{Hom}_{\mathrm{Ab}}(Z, A^*) \simeq \operatorname{Hom}_{k\text{-alg}}(k(Z), A),
\]
\[
  G = \operatorname{Spec} k(Z)
\]
($k(Z)$ : \uncertain{algèbre de groupe}) \add{\uncertain{lin.\ sur}} $M$ [\ill{} \ill{} qu'on entend : sur $W(M) : A \mapsto M \otimes A$, \ill{} au choix pour

\page{149}
\marginal{en marge gauche, écrit en long et relié par une accolade au paragraphe suivant : « La donnée d'une telle $\Delta$ revient à la donnée des $\Delta_n : M_n \to \sum_{i+j=n} M_i \otimes M_j$, i.e.\ de $\prod_{i+j=n} V(M_i) \times V(M_j) \to V(M_{i+j})$ », avec sous $V(M_i) \times V(M_j)$ et $V(M_{i+j})$ les notations $V(M)^i \times V(M)^j$ et « $V(M)^n$ ».}
contragrédiente, sur $V(M)$].

Une structure de coalg.\ $\Delta$ sur $M$ est dite compatible avec la graduation, si
\[
  \Delta : M \longrightarrow M \otimes M
\]
est compatible avec la « graduation totale » ($\Delta(M_n) \subset \sum_{i+j=n} M_i \otimes M_j$), \uncertain{ce qui revient au même}, si la loi de composition sur $V(M)$ est compatible avec les graduations, i.e.\ provient d'homs
\[
  \underset{V(M)^i}{V(M_i)} \times \underset{V(M)^j}{V(M_j)} \longrightarrow \underset{V(M)^{i+j}}{V(M_{i+j})} .
\]
Il revient aussi au même de dire que la multiplication dans $V(M)$ (\uncertain{donc la} comultiplication de $M$) commute à l'action de $G$.
\marginal{NB dans ce pt de vue « familles de modules indexés par $M_i$ », on tient compte que \ill{} $\Delta$ \ill{} revient à se donner des $\Delta^{(ij)} : M_{i+j} \to M_i \otimes M_j$ \ill{} \ill{} \uncertain{loc.\ finis} \ill{} : $V(M)^i \times V(M)^j \to V(M)^{i+j}$ \ill{} \ill{}.}

Bien sûr, les \ill{} (resp.\ de coalgèbres graduées) correspondent aux \ill{} des $\mathcal{O}_k$-Alg.\ correspondantes (éventuellement graduées).

Condition d'anticommutativité, lorsqu'on dispose d'un hom
\[
  \text{\struck{\ill{}}}\ Z \longrightarrow \mathbb{Z}/2\mathbb{Z}
\]
(i.e.\ $\mu_2 \to G$) qui permet de \struck{\ill{}} \add{déduire de la graduation de $M$ une} graduation de type $\mathbb{Z}/2\mathbb{Z}$, \uncertain{idem pour} $V(M)$.

\page{150}
\marginal{en marge gauche, écrit en long : « signifie que le composé $M^2 \xrightarrow{\Delta^{11}} M^1 \otimes M^1 \to \operatorname{Sym}^2 M^1$ est nul (\uncertain{degré impair}) » ; un trait le rattache à la condition $x^2 = 0$ ci-dessous.}
On dit \struck{alors} que $\Delta$ (ou la coalgèbre $M$) est anticommutative si \struck{\ill{}} la multiplication sur \struck{$V(M)$} l'algèbre graduée \struck{$V(M) = V(M)^{\mathrm{pair}}$}
\[
  V(M) = V(M)^{\mathrm{pair}} \times V(M)^{\mathrm{impair}}
\]
l'est, i.e.\ si on a (pour $x, y \in V(M)(A)$ homogènes mod $2$)
\[
  xy = (-1)^{\deg x \deg y}\, yx ,
\]
i.e.\ $xy = yx$ si $x$ ou $y$ \add{est de degré pair} \struck{\ill{}}, \struck{$xy = yx$} $xy = -yx$ i.e.\ $xy + yx = 0$ si l'un et l'autre de degré impair). On laisse au lecteur le soin de \uncertain{diagrammatiser} cette condition.

On dit que $\Delta$ (ou la coalgèbre $M$) est alternée si la multiplication de $V(M)$ l'est, i.e.\ si on a pour \struck{\ill{}} $x, y \in V(M)(A)$, homogènes mod $2$, \struck{\ill{}}
\[
  \begin{cases}
    xy = yx & \text{si } x \text{ ou } y \text{ de degré pair}  \\
    x^2 = 0 & \text{si } x \text{ de degré impair.}
  \end{cases}
\]
\note{la condition $x^2 = 0$ est entourée.}
Comme, si $x$ et $y$ de degré impair, \uncertain{alterné} $\Rightarrow$ $xy + yx = 0$, donc
\[
  x^2 = y^2 = (x+y)^2 = 0 ,
\]
\ill{} donc \add{\ill{}} la coalgèbre \uncertain{associée} \ill{} l'Algèbre est anticomm., donc \ill{} anticommutative. Bien sûr, \uncertain{alternance et} anticommutativité \add{\uncertain{la condition}} \struck{\ill{}} \add{\ill{}} \uncertain{des} \struck{\ill{}} \ill{} coalgèbres se conservent par \uncertain{lim.\ quelc.}

On définit les $k$-\add{co}dérivations d'une coalgèbre, et les $k$-antidérivations (\uncertain{degré}

\page{151}
\uncertain{supposé} de degré impair, i.e.\ \uncertain{échangeant parité} des degrés) pour des coalgèbres \struck{\ill{}} comme correspondant aux dérivations resp.\ antidérivations de la $\mathcal{O}_k$-Algèbre $V(M)$ (\uncertain{soit graduée} mod $2$). On a donc que le crochet de deux dérivations est une dérivation, l'anticrochet de deux antidérivations est une antidérivation. On notera que si $\partial$ (\uncertain{venant} à un gr.\ de degrés quelc.) est un endom.\ de $M$ et de degré $\nu$, alors l'endom.\ correspondant de $V(M)$ est de degré $-\nu$ --- sauf si on change les degrés \add{de signe} \struck{\ill{}} par l'astuce
\[
  \underset{V(M^{-i})}{V(M_i)} = V(M)^i = V(M)_{-i} .
\]
Mais pour les histoires de dérivation \add{et antidérivation} \struck{\ill{}} seul un degré mod $2$ importe, donc il n'y a pas \uncertain{lieu} de distinguer entre degré $i$ et $-i$ d'un endom.\ linéaire.

On traite de façon analogue les « coproduits divisés », \struck{\ill{} « \ill{} » d'une coalgèbre} \add{modules \ill{}, \uncertain{sous-coalgèbre} d'une} coalgèbre \ill{} Pour le faire avec la généralité qui convient (\uncertain{par seulement} dans le cas comm., \ill{} mais aussi le cas anticomm.) il faut d'abord disposer des \uncertain{résultats} dans le cas pas « co », mais ordinaire.

\section{Généralités sur coalg.\ et comodules}

\page{152}
\note{titre encadré en tête de page, de la main de Grothendieck.}
Soit $s$ une $k$-coalg.\ (\struck{\ill{}}, ass., unit.)
$C$ un comodule \add{(à g.)} sur $s$,

donc $V_k(C)$ est un Module \add{(à g.)} sur $V_k(s)$ ; \struck{\ill{}} donc \add{\uncertain{l'Alg.}} $V_k(s)$ opère \add{(à g.)} sur $V_k(C)$, donc l'algèbre $V_k(s)(k) = s^k$ opère \add{(à g.)} sur $V_k(C)$, donc \add{à droite} sur $C$. Donc $C$ est un module \add{(à droite)} ordinaire sur $s^k$, et sa structure de $k$-module \uncertain{s'en déduit} (par $k \to s^k$).

Quand peut-on récupérer la structure de comodule \add{(à g.)} sur $s$, connaissant la structure de module \add{(à dr.)} sur $s^k$ ? [NB. il est clair que si on connaît la structure de \struck{\ill{}} module \add{à dr.}\ de $C \otimes_k k'$ sur $s^{k'}$ $\forall k'$, alors on connaît la structure de comodule de $C$ sur $s$ ; de façon plus précise, une telle structure de comodule \add{à g.}\ revient à \uncertain{la} donnée d'une structure de module, \add{à droite} de
\[
  W(C) = (k' \longmapsto C \otimes_k k') \quad \text{sous} \quad V(s) = (k' \longmapsto s^{k'} = \operatorname{Hom}_{k\text{-mod}}(s, k')) \text{.]}
\]
Soit $u \in s^k$, $u : s \to k$, alors l'op.\ $u_C$ de $u$ sur $C$ est le composé

\begin{tikzcd}
C \arrow[r, "\Delta_C"] \arrow[rr, bend right=30, "u_C"'] & C \otimes s \arrow[r, "\mathrm{id}_C \otimes u"] & C \otimes k \simeq C
\end{tikzcd}

donc la connaissance des $u_C$ équivaut à la connaissance du composé
\[
  C \xrightarrow{\ \Delta_C\ } C \otimes s \xrightarrow{\ \mathrm{id}_C \otimes u\ } \prod_{u \in s^k} C \otimes k
\]
et cela est bien si
\[
  C \otimes s \longrightarrow \prod_{u \in s^k} C \otimes k
\]
est injectif.

\page{154}
Si $C$ est plat sur $k$, et $s \hookrightarrow k^I$ pour $I$ convenable
i.e.\ $s \to k^{(s^k)}$ $\bigl(x \mapsto (u(x))_{u \in s^k}\bigr)$ est injectif,
alors $C \otimes s \to C \otimes_s k^I$ est injectif, mais
\[
  C \otimes_k k^{(I)} \longrightarrow (C \otimes_k k)^I = C^I
\]
\[
  \operatorname{Hom}(k^{(I)}, C)
\]
\note{ce $\operatorname{Hom}$ est écrit sous $(C \otimes_k k)^I$, relié par un signe d'égalité vertical.}
est-il injectif ? \uncertain{Oui}, par passage à la limite inductive sur $I$, on se ramène au cas $I$ fini // \ill{} \uncertain{on conclut} sans hypothèse de platitude sur $C$ en supposant que $s$ est \emph{libre} sur $k$ …

\note{deux traits obliques en marge gauche.}
\[
  s = k[U] , \qquad s^k = k\{\{T\}\} \quad \text{structure de \uncertain{k-alg.\ ordinaire}}
\]
\[
  C = k[U][Y_0, \ldots, Y_n] = k[U, Y_0, \ldots, Y_n] , \qquad C^k = k\{\{T\}\}\{\{X_0, \ldots, X_n\}\}
\]
structure de $k$-alg.\ ordinaire qui est une $s^k$-alg.

$C^k$ opère sur $C$ par les
\[
  T^{(i)} X_0^{(i_0)} \cdots X_n^{(i_n)} \quad \text{opérant par l'op.\ diff.} \quad D_U^{(i)} D_{Y_0}^{(i_0)} \cdots D_{Y_n}^{(i_n)} .
\]
\struck{NB} Donc $s^k$ opère par les $D_U^{(i)}$.

\emph{Remarque.} Toutes ces opérations sont \struck{\uncertain{nilpotentes}} \add{loc.} \uncertain{nilpotentes}, de façon précise si $P \in C$ fixé, $P$ de degré $\leqslant n$, si $f \in C^k$ est d'\emph{ordre} $> n$, alors $D_f P = 0$ \ill{} fonction \ill{} $f \in s^k$ d’ordre > …

\note{deux traits obliques en marge gauche.}
Supposons $s$ $k$-\struck{\ill{}} coalgèbre cocommutative, unitaire, libre comme $k$-module, quelles sont les structures de $s^k$-\uncertain{Algèbres} sur \add{les $k$-modules} $C$ qui proviennent d'une structure de $k$-comodule ? \add{\uncertain{i.e.} comme $s$ libre -- cf.\ ci-dessus}
\[
  C \longrightarrow C \otimes_k s \hookrightarrow \operatorname{Hom}_k(s^k, C)
\]
\[
  \text{\struck{$C \otimes_k \operatorname{Hom}_k((s^k)^k, k)$}} \longrightarrow \operatorname{Hom}_k(s^k, C \otimes k)
\]
\note{la dernière ligne est en partie biffée ; sous l'argument biffé on lit $(s^k)^k$. Le passage reste en suspens en bas de page.}

\page{156}
On voit que \uncertain{dans} OK \uncertain{ces} $\forall x \in C$, l'application
\[
  i_x : u \longmapsto u_C\, x , \qquad s^k \longrightarrow C
\]
est dans $C \otimes_k s$, \struck{\ill{}} i.e.\ \uncertain{en termes} d'une base $(e_\alpha)_{\alpha \in A}$ de $s$) on a une formule
\[
  u_C\, x = \Bigl(\underbrace{\sum x_\alpha \otimes e_\alpha}_{\substack{\text{somme finie,} \\ x_\alpha \in C}}\Bigr)(u) = \sum u(e_\alpha)\, x_\alpha
\]
\uncertain{en posant}, on \uncertain{pose}
\[
  u(e_\alpha) = u_\alpha \quad \text{dans} \quad u \longmapsto (u_\alpha)_{\alpha \in A} ,\ s^k \simeq k^A ,
\]
\[
  u . x = \sum u_\alpha x_\alpha .
\]
\note{triple trait vertical en marge gauche des trois lignes qui suivent.}
Cela signifie aussi que $u \mapsto u.x$ est continue pour \uncertain{la} top.\ discrète sur $C$ et la top.\ de convergence simple sur $s^k$. On en conclut que
\[
  s^k \longrightarrow \operatorname{Hom}_k(C, C)
\]
est continu, quand on met dans les deux cas la top.\ de conv.\ simple-discrète …

Supposons maintenant que $s$ soit graduée \add{$= \sum s_n$} (\uncertain{graduée} $\Delta$) donc $s^k$ gradué,
\[
  s^k = \prod_{n \in \mathbb{Z}} {s_n}^k = \prod_{n \in \mathbb{Z}} (s^k)_n .
\]
\note{sous le dernier membre, une accolade porte « déf ».}
Soit $C$ \add{$= \sum C_m$} \uncertain{comodule} gradué --- donc …
$C^k = \prod (C_m)^k = \prod (C^k)_m$ \struck{comodule} gradué sur $s^k = \prod (s^k)_n = \prod ({s_n}^k)$.
Alors \struck{si $u \in s^k$} si $u \in (s^k)_n$, $(u_C)^k : V(C_m) \to V(C_{m+n})$,
donc $u_C : C_{m+n} \to C_m$ i.e.\ $u_C$ \emph{diminue} les degrés de $n$ ! (Pour qu'il les augmente il faudrait changer de signe les degrés de $C$ \struck{\ill{}} \add{ou} plutôt exiger que $V(C)^m \overset{\mathrm{def}}{=} V(C_{-m})$ …)

\page{158}
Donc la catégorie des $s$-comodules gradués équivaut à la catégorie des $s^k$-modules gradués $C$ \struck{\ill{}} qui sont « \struck{\ill{}} \uncertain{loc.\ nilpotents} » i.e.\ tels que $\forall x \in C$, $u \mapsto u.x$ de $s^k$ dans $C$ soit continue pour la top.\ simple sur $s^k$ (\struck{$=$} \uncertain{top.\ produit} sur $k^A$).

\page{159}
\note{notes de sa main dans la marge d'un feuillet dactylographié retourné ; elles reprennent le calcul de la p. 156.}
\[
  \text{\struck{\ill{}}}\ (s^k)^n = (s_{-n})^k
\]
\[
  u : \underset{V(C_{-m})}{V(C_m)^m} \longrightarrow \underset{V(C_{-m-n})}{V(C)^{m+n}}
\]
\note{la suite de ces notes manque dans le lot ; l'argument se poursuit au-delà de la p. 160.}

\end{document}
